% Korean edition of sections/s2_1_transport.tex.
\seclabel
\claim{\secnum{2.1}{sec:transport}\quad HTTP/2 연결 하나를 재사용하면 요청 시간이 106.5\ms{}에서 16.0\ms{}로 준다}

\lead{사이트는 TLS 1.3 위에서 HTTP/2를 제공하며 HTTP/3는 알리지 않는다.} TLS 핸드셰이크는 ALPN으로 HTTP/2를 고르고, 암호 스위트는 \code{TLS_AES_256_GCM_SHA384}, 인증서 체인은 Let's Encrypt의 인증서 4개다. 응답에는 서버가 HTTP/3(QUIC)를 알릴 때 쓰는 \code{alt-svc} 헤더가 없다. HTTP/3가 없다는 판단은 이 헤더가 없다는 사실에만 근거하며, UDP 443 포트는 확인하지 않았다.

\lead{nginx 프런트엔드가 요청을 ASP.NET 애플리케이션으로 넘긴다.} 요청 본문이 잘못되면 nginx 고유의 HTML 오류 페이지가 나온다(\ref{sec:auth}절). 애플리케이션의 오류 메시지와 SignalR 허브로 프레임워크를 식별했다. \ref{sec:auth}절과 \ref{sec:feed}절에서 인용하는 ASP.NET 기본값은 이 식별에 근거한다.

\lead{REST 요청은 HTTP/2 연결 하나를 공유하고, 실시간 피드는 두 번째 연결이 필요하다(\figref{layers}).} REST 요청 하나는 \code{Authorization: Bearer} 헤더를 실은 HTTP/2 스트림 하나로 전달된다. 이후 요청은 같은 연결에 새 스트림을 열어 DNS, TCP, TLS를 건너뛴다. 실시간 피드는 협상(negotiate)용 HTTP/2 POST로 시작한 뒤, 두 번째 연결을 열고 HTTP/1.1 \code{Upgrade: websocket} 요청으로 업그레이드한다. 이 두 번째 연결은 서버가 강제한다. 클라이언트가 HTTP/2 위로 WebSocket을 실으려면 서버가 먼저 \code{ENABLE_CONNECT_PROTOCOL} 설정을 보내야 하는데[1], 서버의 HTTP/2 SETTINGS 프레임에 이 설정이 없다. 따라서 브라우저를 포함한 모든 클라이언트가 피드용 연결을 따로 연다.

% Korean edition of figs/layers.tex.
\fig[wide]{layers}{REST 요청은 HTTP/2 연결 하나를 공유하고, 실시간 피드는 두 번째 WebSocket 연결이 필요하다}[채널별 클라이언트 측 프로토콜 스택: 애플리케이션부터 연결 설정까지]{%
\begin{tikzpicture}[x=1mm,y=1mm,
  box/.style={draw=navy,line width=0.5pt,fill=white,minimum width=30mm,minimum height=8mm,inner sep=0pt,font=\fontsize{7.5}{9}\selectfont,anchor=north west},
  lname/.style={anchor=west,align=left,font=\fontsize{7.5}{9}\selectfont\bfseries\color{navy}},
  note/.style={anchor=west,align=left,text width=69mm,inner sep=0pt,font=\fontsize{7}{8.4}\selectfont\color{greyD}},
  hd/.style={anchor=south,align=center,font=\fontsize{7.5}{9}\selectfont\bfseries\color{navy}}]
% column heads
\node[hd] at (58,-1.2) {REST\\[-0.5mm]{\fontsize{7}{8.4}\selectfont\mdseries\color{greyD}\code{/api}}};
\node[hd] at (92,-1.2) {실시간 피드\\[-0.5mm]{\fontsize{7}{8.4}\selectfont\mdseries\color{greyD}\code{/hubs/main}}};
% stacks: rows of 8 mm from y = -2
\foreach \y/\n/\a/\b in {-2/애플리케이션/{REST API}/SignalR, -10/프레이밍/{HTTP/2 스트림}/WebSocket, -18/HTTP/{HTTP/2}/{HTTP/1.1}, -26/TLS/{TLS 1.3}/{TLS 1.3}, -34/TCP/TCP/TCP}{
  \node[box] at (43,\y) {\a};
  \node[box] at (77,\y) {\b};
  \node[lname] at (23,\y-4) {\n};
}
% DNS, before connecting
\node[box,draw=greyD,dashed,minimum width=64mm] at (43,-45) {DNS};
\node[lname,text=greyD] at (23,-49) {연결 전};
% bracket over TLS, TCP and DNS of both stacks
\draw[navy,line width=0.6pt] (21.5,-26.3) -- (20,-26.3) -- (20,-52.7) -- (21.5,-52.7);
\draw[navy,line width=0.6pt] (20,-39.5) -- (18.5,-39.5);
\node[anchor=east,align=right,font=\fontsize{7}{8.4}\selectfont\bfseries\color{navy}] at (18,-39.5) {새 연결마다\\다시 치르는\\비용};
% short notes, one per row
\node[note] at (111,-6)  {REST: 요청마다 Bearer 토큰.\\피드: 연결을 열 때 토큰을 한 번 검사};
\node[note] at (111,-14) {REST: 요청마다 스트림 하나, 연결 하나에 여러 개.\\피드: JSON 레코드, 각각 0x1E로 끝남};
\node[note] at (111,-22) {REST: ALPN으로 선택. 피드: 협상용 HTTP/2 POST 후,\\별도 연결에서 \code{Upgrade: websocket}};
\node[note] at (111,-30) {왕복 1회. 피드는 자체 핸드셰이크를 따로 함};
\node[note] at (111,-38) {REST: 15\textasciitilde 30\,s 유휴 후 서버가 닫음.\\피드: 3\,s마다 요청해 열어 둠};
\node[note] at (111,-49) {\code{contest.timefolio.net} $\rightarrow$ 221.168.38.40};
\end{tikzpicture}
}{상자 하나가 프로토콜 하나다. 두 채널 모두 같은 DNS 조회에서 시작한다. 이후 REST 요청은 같은 연결에 새 HTTP/2 스트림을 열어 DNS, TCP, TLS를 건너뛴다. 서버의 HTTP/2 SETTINGS 프레임에 \code{ENABLE_CONNECT_PROTOCOL}이 없어 브라우저를 포함한 어떤 클라이언트도 HTTP/2 위의 WebSocket[1]을 쓸 수 없으므로, 실시간 피드에는 두 번째 연결이 필요하다. 서버에서는 nginx 프런트엔드가 요청을 ASP.NET 애플리케이션으로 넘긴다. SETTINGS 프레임은 Python \code{h2} 라이브러리로 읽었다. 프런트엔드는 오류 페이지로, 애플리케이션은 오류 메시지로 식별했다.}{DART Quant Team}


\lead{새 연결은 106.5\ms{}, 재사용 연결은 16.0\ms{}가 걸리므로 첫 요청에서 연결 설정에 약 90\ms{}가 든다.} 같은 엔드포인트(\code{Portfolio/Violations})로 curl 한 번 실행에 요청 5개를 연달아 보냈고(1회에 약 170\ms), 이를 1\,s 간격으로 10회 반복했다. 각 회의 첫 요청은 DNS 조회, TCP 연결, TLS 핸드셰이크 비용을 치르고 나머지 4개는 연결을 재사용한다(\figref{timing}). 단계별 중앙값의 합은 106.0\ms{}이고, 합계의 중앙값 106.5\ms{}는 다른 회에서 나온다.

% Korean edition of figs/timing.tex.

\fig[wide]{timing}{새 연결은 106.5\ms, 재사용 연결은 16.0\ms{}가 든다}[\code{Portfolio/Violations}에 보낸 curl 요청 하나, 10회의 중앙값(재사용 요청 40개). 시간은 아래로 흐르며 두 패널의 축척은 같다]{%
\begin{tikzpicture}[x=1mm,y=1mm,
  f/.style={font=\fontsize{7}{8.2}\selectfont},
  host/.style={font=\fontsize{7}{8.2}\selectfont\bfseries,text=navy,anchor=south,align=center},
  tl/.style={greyM,line width=0.6pt,dotted,-{Stealth[length=1.4mm]}},
  msg/.style={-{Stealth[length=1.6mm]},line width=0.7pt,navy},
  rep/.style={-{Stealth[length=1.6mm]},line width=0.7pt,navyM},
  ml/.style={f,text=navy,sloped,above,inner sep=0.6pt},
  stg/.style={line width=0.5pt,navy},
  sl/.style={f,anchor=east,align=right,text=navy},
  ptitle/.style={font=\fontsize{7.5}{9}\selectfont\bfseries,text=navy,anchor=west}]
\def\tT#1{-0.72*(#1)}% time in ms -> y in mm (0.72 mm per ms, both panels)
% ============================================ panel (a): new connection
\def\Rx{0}\def\Cx{22}\def\Sx{58}
\node[ptitle] at (-24,14) {(a) 새 연결: 요청 전에 DNS, TCP, TLS를 거친다};
\node[host] at (\Rx,2.5) {DNS\\리졸버};
\node[host] at (\Cx,2.5) {클라이언트\\(curl)};
\node[host] at (\Sx,2.5) {서버\\(nginx)};
\foreach \x in {\Rx,\Cx,\Sx} {\draw[tl] (\x,1) -- (\x,{\tT{110}});}
% DNS: the split inside the 49.3 ms is not measured; arrows drawn one-way 3 ms
\draw[msg] (\Cx,{\tT{0}}) -- node[ml]{A 질의} (\Rx,{\tT{3}});
\fill[navyXL] (\Rx-1,{\tT{3}}) rectangle (\Rx+1,{\tT{46.3}});
\draw[rep] (\Rx,{\tT{46.3}}) -- node[ml]{A 응답} (\Cx,{\tT{49.3}});
% TCP: one RTT
\draw[msg] (\Cx,{\tT{49.3}}) -- node[ml]{SYN} (\Sx,{\tT{54.85}});
\draw[rep] (\Sx,{\tT{54.85}}) -- node[ml,pos=0.78,fill=white,inner sep=0.6pt]{SYN-ACK} (\Cx,{\tT{60.4}});
% TLS 1.3: one RTT, then the rest of curl's TLS time on the client
\draw[msg] (\Cx,{\tT{60.4}}) -- node[ml,pos=0.66]{ACK, ClientHello} (\Sx,{\tT{65.95}});
\draw[rep] (\Sx,{\tT{65.95}}) -- node[ml,below,fill=white,inner sep=0.6pt,pos=0.35]{ServerHello ... Finished} (\Cx,{\tT{71.5}});
\fill[navyXL] (\Cx-1,{\tT{71.5}}) rectangle (\Cx+1,{\tT{89.3}});
\node[f,anchor=west,text=greyD,align=left] at (\Cx+1.8,{\tT{81.4}}) {curl TLS 시간의 나머지};
% request to first byte
\draw[msg] (\Cx,{\tT{89.3}}) -- node[ml]{Finished, \code{HEADERS} (요청)} (\Sx,{\tT{94.85}});
\fill[navyXL] (\Sx-1,{\tT{94.85}}) rectangle (\Sx+1,{\tT{100.45}});
\node[f,anchor=west,text=greyD] at (\Sx+1.8,{\tT{97.65}}) {서버};
\draw[rep] (\Sx,{\tT{100.45}}) -- node[ml,below,inner sep=0.4pt]{\code{HEADERS} 200, \code{DATA}} (\Cx,{\tT{106}});
% stage brackets (left)
\draw[stg] (-3+1,{\tT{0}-0.2}) -- ++(-1,0) -- (-3,{\tT{49.3}+0.2}) -- ++(1,0);
\node[sl] at (-4.6,{\tT{24.65}}) {DNS 조회\\\textbf{49.3}\ms\\{\color{greyD}WSL에}\\{\color{greyD}DNS 캐시 없음}};
\draw[stg] (-3+1,{\tT{49.3}-0.2}) -- ++(-1,0) -- (-3,{\tT{60.4}+0.2}) -- ++(1,0);
\node[sl] at (-4.6,{\tT{54.85}}) {TCP \textbf{11.1}\ms\\= 왕복 1회};
\draw[stg] (-3+1,{\tT{60.4}-0.2}) -- ++(-1,0) -- (-3,{\tT{89.3}+0.2}) -- ++(1,0);
\node[sl] at (-4.6,{\tT{74.85}}) {TLS 1.3 \textbf{28.9}\ms\\{\color{greyD}핸드셰이크 왕복 1회,}\\{\color{greyD}Python에서 12.4\ms}};
\draw[stg] (-3+1,{\tT{89.3}-0.2}) -- ++(-1,0) -- (-3,{\tT{106}+0.2}) -- ++(1,0);
\node[sl] at (-4.6,{\tT{97.65}}) {요청에서 첫 바이트까지\\\textbf{16.7}\ms};
\draw[stg] (\Sx+12,{\tT{0}}) -- ++(1,0) -- (\Sx+13,{\tT{106}}) -- ++(-1,0);
\node[f,text=navy,rotate=-90,anchor=south] at (\Sx+14.5,{\tT{53}}) {합계 \textbf{106.5}\ms{} (합계의 중앙값, 단계별 중앙값의 합은 106.0\ms)};
% ============================================ panel (b): reused connection
\def\Cb{94}\def\Sb{126}
\node[ptitle] at (\Cb-8,14) {(b) 재사용 연결: 요청만 보낸다};
\node[host] at (\Cb,2.5) {클라이언트\\(curl)};
\node[host] at (\Sb,2.5) {서버\\(nginx)};
\foreach \x in {\Cb,\Sb} {\draw[tl] (\x,1) -- (\x,{\tT{24}});}
\draw[msg] (\Cb,{\tT{0}}) -- node[ml]{\code{HEADERS} (요청)} (\Sb,{\tT{5.55}});
\fill[navyXL] (\Sb-1,{\tT{5.55}}) rectangle (\Sb+1,{\tT{10.45}});
\draw[rep] (\Sb,{\tT{10.45}}) -- node[ml,below,inner sep=0.4pt]{\code{HEADERS} 200, \code{DATA}} (\Cb,{\tT{16}});
\draw[stg] (\Cb-3+1,{\tT{0}-0.2}) -- ++(-1,0) -- (\Cb-3,{\tT{16}+0.2}) -- ++(1,0);
\node[sl] at (\Cb-4.6,{\tT{8}}) {\textbf{16.0}\ms};
\node[f,anchor=north west,align=left,text=navy] at (\Cb-6,{\tT{24}-2}) {왕복 1회(11.1\ms)와 서버에서 약 5\ms:\\열린 연결 위의 새 HTTP/2 스트림};
% run-to-run spread strip: min to max over 10 runs (40 for reused), tick = median; 0.18 mm per ms
\begin{scope}[shift={(\Cb+14,{\tT{24}-17})}]
\def\s{0.16}
\node[f,anchor=south west,font=\fontsize{7}{8.2}\selectfont\bfseries,text=navy] at (-20,1.6) {회차별 분포, 최솟값\textasciitilde 최댓값(ms)};
\foreach \lab/\lo/\md/\hi/\r in {DNS/36.6/49.3/94.1/0, TCP/9.4/11.1/12.3/1, TLS/17.8/28.9/40.4/2, 첫 바이트까지/15.7/16.7/19.8/3, {합계, 새 연결}/80.5/106.5/156.9/4, {합계, 재사용}/14.8/16.0/20.3/5} {
  \pgfmathsetmacro{\yy}{-3.9*\r}
  \node[f,anchor=east,text=navy] at (-0.8,\yy) {\lab};
  \draw[navyL,line width=1.8pt] ({\s*\lo},\yy) -- ({\s*\hi},\yy);
  \draw[navy,line width=1pt] ({\s*\md},\yy-1.3) -- ({\s*\md},\yy+1.3);
  \node[f,anchor=west,text=greyD] at ({\s*160+0.8},\yy) {\lo{}\textasciitilde\hi};
}
\draw[greyM,line width=0.4pt] (0,-21.5) -- ({\s*160},-21.5);
\foreach \t in {0,50,100,150} {\draw[greyM,line width=0.4pt] ({\s*\t},-21.5) -- ({\s*\t},-22.3); \node[f,anchor=north,text=greyD] at ({\s*\t},-22.3) {\t};}
\end{scope}
% legend and scale bar
\begin{scope}[shift={(-24,{\tT{110}-6})}]
\draw[msg] (0,0) -- (7,0); \node[f,anchor=west] at (8,0) {클라이언트에서 서버로};
\draw[rep] (32,0) -- (39,0); \node[f,anchor=west] at (40,0) {응답};
\fill[navyXL] (54,-1.2) rectangle (56,1.2); \node[f,anchor=west] at (57,0) {로컬 작업};
\draw[greyD,line width=0.6pt] (78,{-0.72*5}) -- (78,{0.72*5});
\foreach \t in {-5,5} {\draw[greyD,line width=0.6pt] (77.2,{0.72*\t}) -- (78.8,{0.72*\t});}
\node[f,anchor=west,text=greyD] at (79.5,0) {두 패널 모두 10\ms};
\end{scope}
\end{tikzpicture}%
}{편도 시간은 TCP 왕복(RTT) 11.1\ms{}의 절반으로 그렸다. 측정한 것은 단계 경계뿐이다. curl의 \code{time_*} 값은 누적값이며, 단계는 인접한 값의 차이다. ‘요청에서 첫 바이트까지’는 curl TLS 단계의 끝에서 응답 첫 바이트까지다. DNS 조회와 curl TLS 단계 내부의 세부 구간은 측정하지 않았다. 라벨은 중앙값이고, 띠는 10회(재사용 요청 40개)의 범위이며 눈금이 중앙값이다. 요청 50개 모두 HTTP 200을 받았다. 2026년 11월 24일 06:03 KST. 시간-공간 그림은 Kurose와 Ross[6]의 형식을 따랐다.}{DART Quant Team}

% Korean edition of figs/pytimings.tex.
\fig[wide]{pytimings}{Python 클라이언트는 프로세스당 한 번 1.2\textasciitilde 1.3\,s, 재사용 요청마다 31\textasciitilde 39\ms{}를 쓴다}[클라이언트, 시점, 연결 상태별 REST 요청 하나의 시간, 2026년 11월 24일]{%
{\tabfont\begin{tabularx}{\linewidth}{>{\RaggedRight\arraybackslash}p{40mm} >{\RaggedRight\arraybackslash}p{42mm} R{17mm} @{\hspace{6pt}}>{\RaggedRight\arraybackslash}X}
\thd{클라이언트, 시점} & \thd{요청} & \thd{시간} & \thd{시간이 드는 곳}\\
curl, 06:03 (개장 전) & 재사용 연결 & 16.0\ms & 왕복 1회(11.1\ms)와 서버에서 약 5\ms{} (\figref{timing})\\\rowsep
\hl Python, 장중 & 재사용 연결 & 31\textasciitilde 39\ms & curl과의 차이 원인은 분리하지 않음: 클라이언트, 시각, 장중 부하\\\rowsep
Python, 새 프로세스 & 첫 요청 & 1.2\textasciitilde 1.3\,s & 모듈 로딩과 SSL 초기화, 프로세스당 한 번. 가상환경이 Windows 드라이브에 있음\\\rowsep
Python, 같은 프로세스 & 두 번째 클라이언트의 첫 요청 & 130\ms & 새 연결 비용만. 일회성 비용은 이미 치름\\
\end{tabularx}}}{음영 행은 시작 이후 래퍼 요청마다 드는 비용이다. 1.2\textasciitilde 1.3\,s는 시작할 때 첫 주문 전에 \code{warm_up()}이 한 번 치른다. Python 시간은 \code{transport} 실서버 검사(\figref{tests})에서 나왔다.}{DART Quant Team}


\lead{\figref{timing}은 유효한 토큰으로 다시 측정한 결과만 쓴다. 만료된 토큰으로 한 이전 측정은 재사용 요청이 12.2\ms{}까지 빨라서 제외했다.} 그 측정은 상태 코드를 기록하지 않았고, 12.2\ms{}는 데이터에 손대기 전에 토큰을 거부하는 서버와 맞아떨어진다. 재측정에서는 모든 요청의 상태 코드를 기록했다.


\lead{재사용 요청 비용은 네트워크 왕복 1회(11.1\ms)에 양 끝에서 쓰는 약 5\ms{}를 더한 값이다.} curl의 TCP 단계는 SYN 세그먼트를 보낸 때부터 SYN-ACK을 받을 때까지로, 클라이언트와 서버 사이 왕복 1회에 해당한다. 중앙값은 11.1\ms{}, 10회의 표준편차는 0.9\ms{}로 네트워크 경로가 안정적이었다. 재사용 연결의 요청도 요청을 보내고 응답을 받는 왕복 1회가 필요하다. 따라서 16.0\ms{} 중 11.1\ms{}는 네트워크 시간이고, 약 5\ms{}는 서버와 클라이언트가 함께 쓴 시간이다. 둘은 분리하지 않았다. 클라이언트 코드를 빠르게 해도 이 5\ms{}의 일부만 줄일 수 있다.

\lead{새 연결의 두 단계는 서버보다 우리 클라이언트 환경을 더 반영한다.} WSL에서는 curl 측정 때 새 연결마다 DNS 조회에 37\textasciitilde 94\ms{}(중앙값 49.3\ms)가 들었고, 별도의 Python 측정에서는 조회당 30\textasciitilde 39\ms{}가 들었다. 1\,s 간격의 조회가 빨라지지 않았으므로 경로상 어떤 캐시도 응답하지 않았다. 캐싱 리졸버가 있는 시스템이라면 조회 비용을 한 번만 낼 것이다. 그래도 래퍼는 호스트 이름으로 연결한다. 인증서가 \code{contest.timefolio.net}을 이름으로 담고 있어서, IP 주소로 연결하면 조회는 아끼지만 호스트 이름 검증을 잃는다. curl의 TLS 28.9\ms{}에는 curl 자체의 TLS 초기화도 들어 있다. 열린 TCP 연결 위에서 Python으로 잰 TLS 1.3 핸드셰이크 자체는 왕복 1회인 12.4\ms{}였다.

\lead{Python 클라이언트에서는 프로세스의 첫 요청이 1.2\textasciitilde 1.3\,s, 재사용 요청이 31\textasciitilde 39\ms{}였다(\figref{pytimings}).} 장중에 잰 재사용 값은 개장 전에 잰 curl의 16.0\ms{}의 약 두 배다. 두 측정은 클라이언트, 시각, 장중 부하가 모두 다르며 이 원인들을 분리하지 않았다. 1.2\textasciitilde 1.3\,s는 거의 전부 일회성 모듈 로딩과 SSL 초기화다(가상환경이 Windows 드라이브에 있다). 같은 프로세스에서 만든 두 번째 클라이언트는 새 연결에 130\ms{}만 썼다. 따라서 일회성 비용은 프로세스에 딸려 오고, Python에서 새 연결 비용은 약 130\ms{}다. 래퍼는 시작할 때 \code{warm_up()}을 호출해 \code{Portfolio/Session}(주문 가능 시간, \ref{sec:queries}절)에 가벼운 GET 하나를 보내고 그 401은 무시한다. 그래서 일회성 비용을 첫 주문 전에 치른다(D18).

\lead{간격마다 한 번씩 시험한 결과, 서버의 유휴 연결 종료 시점은 15\textasciitilde 30\,s 사이다.} httpx는 기본으로 5\,s 유휴 후 연결을 끊기 때문에[4] 첫 시도에서는 서버의 동작이 가려졌다. 이 한도를 600\,s로 올리자 5\,s, 15\,s 침묵 뒤에는 연결이 재사용됐고 30\,s 뒤에는 새 연결로 바뀌었다. 간격마다 한 번만 시험했으므로 종료 시점은 연결마다 다를 수 있다. HTTP/2 서버는 보통 GOAWAY 프레임으로 유휴 종료를 알리고, 클라이언트는 다음 요청 때 새 연결을 연다. 프레임은 캡처하지 않았으므로 이 메커니즘은 관찰이 아니라 가정이다. GOAWAY가 오는 도중에 요청을 보낼 때만 요청이 닫힌 소켓을 만난다. 이 경합에 걸린 GET은 한 번 더 보내지만, 서버가 받았을 수 있는 주문은 다시 보내지 않으므로(\ref{sec:orders}절) 래퍼가 서버보다 먼저 유휴 연결을 닫는다.

\verdict{래퍼는 프로세스가 살아 있는 동안 \code{httpx.AsyncClient(http2=True)} 하나를 유지한다. \code{keepalive_expiry}는 10\,s로 둔다. 연결이 여전히 재사용된 가장 긴 유휴 시간인 15\,s보다 짧으므로, 서버보다 래퍼가 먼저 유휴 연결을 닫는다(D1, D2). 직전 요청 후 10\,s 안의 요청은 연결을 재사용하고, 그보다 긴 간격 뒤에는 새 연결을 연다(Python에서 약 130\ms). 주문 직전에는 사전 검사의 \code{ValidateProduct} GET(\ref{sec:guidelines}절)이 연결을 다시 열어 두므로, \code{AddOrder} 자체는 열린 연결로 나간다.}
